\section{\bench and Evaluation Protocol}
\label{app:benchmark}
\suppressfloats[t]

\subsection{Benchmark Design}
\label{app:tasks}
\bench comprises fourteen tasks in five families (Fig.~\mainref{fig:bench_overview}):
grasping a moving target, placing an object into a moving receptacle, passing a wand
along a rotating wire, capturing rolling or sliding objects, and grasping after an
impact. Together, they cover translation, rotation, temporary occlusion, constrained
passage, and abrupt changes of the target's motion. We will release the tasks, the
accepted scenarios, the success checks, and the demonstrations.

\paragraph{Environment.} The tasks run in MuJoCo 3.11.0 and robosuite 1.5.2
\cite{mujoco,robosuite} with a Franka Panda arm and a parallel gripper, driven by
robosuite's operational-space pose controller at a control period of 50\,ms and a
physics step of 2\,ms. A fixed third-person agent-view camera at
$(0.68,-0.58,1.62)$\,m with a $40^\circ$ vertical field of view renders
$224\times224$ RGB-D images. An action is an end-effector increment in the robot
base frame, a translation and a rotation vector, together with a gripper command; the
shared action adapter clips each command to 0.10\,m per translation axis and 0.5\,rad
per rotation-vector component.

\paragraph{Success and termination.} Table~\ref{tab:tasks} lists the ordered success
conditions. We check them at every physics step, an interrupted continuous condition
resets its counter, and an episode ends at its first terminal event: success, an
invalid simulator state, a boundary event (leaving the workspace, a drop, or an
expired goal window), a safety event such as unsafe contact, a task violation, or the
time limit.

\begin{table*}[!t]
\centering\fontsize{9}{11}\selectfont
\setlength{\tabcolsep}{4pt}
\renewcommand{\arraystretch}{1.13}
\caption{Native success conditions in the reported fourteen-task suite.
Thresholds come from the evaluator source and task configurations bound to the
accepted FUTURE0, BASE2, and HISTORY1 packages. A task must pass its preceding
stages and avoid terminal failures. All tasks run at 20\,Hz with a 2\,ms physics
step; durations below refer to physical time. COM denotes center of mass.}
\label{tab:tasks}
\begin{tabularx}{\textwidth}{@{}p{0.19\textwidth}Xp{0.07\textwidth}@{}}
\toprule
Task & Ordered success conditions & Limit \\
\midrule
grasp;\newline t1\_belt, t1\_xyz,\newline t1\_so3, t1\_occ
& Approach within 0.10\,m; bilateral fingerpad contact for 0.05\,s; collision-geometry
bottom at least 0.05\,m above support; retain contact and clearance for 0.50\,s.
T1 adds conveyor drive, translation, rotation, or a physical occluder.
SO3 has no additional success-angle threshold; unsafe occluder contact terminates OCC.
& 160 steps\newline 8\,s \\
t2\_trans
& Approach within 0.10\,m. Entire object collision geometry inside the receptacle:
planar clearance and floor gap both $\geq-0.002$\,m, top no higher than the opening.
Release requires an opening command $<-0.5$ and no fingerpad contact.
Relative linear/angular speeds $\leq0.08$\,m/s and $0.8$\,rad/s for 0.10\,s,
followed by a 0.50\,s stable dwell.
& 200\newline 10\,s \\
t2\_se3
& Approach within 0.05\,m and 0.12\,rad. Full T-shaped containment:
planar clearance $\geq0.002$\,m, floor gap in $[-0.002,0.002]$\,m,
tilt $\leq0.05$\,rad, top below opening.
After release, relative speeds $\leq0.02$\,m/s and $0.04$\,rad/s
for 0.10\,s, then a 0.50\,s dwell. The 0.12\,rad gate applies to approach.
& 200\newline 10\,s \\
t3\_phase
& Start with the native pre-grasp. Phase alignment: distance $\leq0.090$\,m and
loop-direction error $\leq0.20$\,rad. Enter within 0.030\,m at progress $\leq0.15$;
reach historical maximum progress $\geq0.60$, then $\geq0.85$ and endpoint
distance $\leq0.055$\,m.
& 420\newline 21\,s \\
t3\_full
& Acquire the wand with bilateral contact for 0.05\,s, then complete the phase
stages above. In a later sample, retain current progress $\geq0.85$ and move
at least 0.045\,m from both path and endpoint. Finally open ($<-0.5$)
with no fingerpad contact. Releasing before exit fails.
& 640\newline 32\,s \\
t4\_sphere
& Approach within 0.10\,m; first qualifying contact after ramp exit; any tray contact before success
fails. Bilateral hold for nominal 25\,ms (12 physics samples, 24\,ms).
Raise COM by 0.05\,m from its height when the bilateral-hold criterion is first met, and retain
contact and height for 0.50\,s.
& 120\newline 6\,s \\
t4\_sphere\_tabletop;\newline t4\_slider
& Approach within 0.10\,m; contact after ramp exit, while supported, in the capture region, and moving
at planar speed $\geq0.08$\,m/s. Bilateral hold for 0.05\,s, COM lift
$\geq0.05$\,m from the first attainment of that hold criterion, then contact and height for 0.50\,s.
& 120\newline 6\,s \\
t4\_can
& Same supported-capture sequence. Before first fingerpad contact, a can-axis
absolute dot product with vertical $>0.45$ fails as unstable tipping.
The configuration does not require a transverse closing axis.
& 120\newline 6\,s \\
t5\_replan
& Exactly one selected-side impact group completes before robot contact.
At 30\,ms after the physical impact-end timestamp, branch-signed lateral speed $\geq0.035$\,m/s and yaw rate
$\geq0.35$\,rad/s. Approach within 0.10\,m; bilateral hold with grasp-axis
error $\leq18^\circ$ for 0.05\,s; COM lift $\geq0.05$\,m from the first
attainment of that hold criterion and stable hold
for 0.50\,s. Impact-group lag is 12\,ms and maximum contact is 0.10\,s.
& 120\newline 6\,s \\
\bottomrule
\end{tabularx}
\end{table*}


\paragraph{Roles.}
\label{app:roles}
Table~\ref{tab:roles} assigns the bodies of each task to the target and context roles
of Sec.~\mainref{sec:formulation}. A moving receptacle is a target, since its pose
decides containment. The wire plays both roles: its pose at the entry is a target,
which the wand must match within 30\,mm, and its geometry along the loop is context,
which any contact violates. Conveyors, ramps, and strikers are drivers. The table also
lists the error terms of Eq.~\mainref{eq:budget} that each task exercises
(Appendix~\ref{app:analysis}) and the two task groups of Fig.~\mainref{fig:analysis},
which follow from these terms and were fixed before the ablations: tasks whose terms
include $r\theta$ are rotation-sensitive, and tasks with an obstacle clearance or a
driver are context-active.

% Protocol table: interaction roles, tolerances, speeds, critical leads, and exercised budget terms per Dynamic-14 task (from the task definitions and sampling ranges; bracketed values are computed from the task geometry before testing).
\begin{table*}[t]
\centering\footnotesize
\caption{\tablelang{\textbf{Roles and tolerances of the \bench tasks.} The target is the body whose pose error
enters a contact or containment tolerance $\varepsilon$. The context comprises obstacles, avoided with clearance $m$, and drivers that move the target.
Terms: error terms that a task exercises (Eq.~\mainref{eq:budget} and Appendix~\ref{app:analysis}). $v$: moving target. $r\theta$: orientation enters the
contact. $e_{\rm dyn}$: non-constant motion (harmonic, gravity, friction, impact). $m$: a context body must be avoided. Occ.:
history gaps. Delay: command delay built into the task. Groups of Fig.~\mainref{fig:analysis}: R, rotation-sensitive; C, context-active.
$^\dagger$Nominal clearance between the wire and the wand's loop for a centered, perpendicular crossing.}{\textbf{\bench 各任务的角色与容差。}目标是位姿误差进入接触或容纳容差 $\varepsilon$ 的物体；环境包括须以间隙 $m$ 避开的障碍物，以及推动目标运动的驱动物。
误差项：任务涉及的误差分解项（式~\mainref{eq:budget} 与附录~\ref{app:analysis}）。$v$：目标运动。$r\theta$：朝向影响接触。$e_{\rm dyn}$：非匀速运动（谐波、重力、摩擦、撞击）。
$m$：须避开某个环境物体。occ.：历史缺失。delay：任务内置的指令延迟。图~\mainref{fig:analysis} 的任务组：R 为旋转敏感组，C 为环境活跃组。
$^\dagger$导线居中且垂直穿过杆端环时两者之间的名义间隙。}}
\label{tab:roles}
\setlength{\tabcolsep}{3.2pt}
\begin{tabularx}{\textwidth}{@{}lll>{\raggedright\arraybackslash}Xccll@{}}
\toprule
\tablelang{Task}{任务} & \tablelang{Code}{代号} & \tablelang{Target}{目标} & \tablelang{Context}{环境} & $\varepsilon$ (mm) & $m$ (mm) & \tablelang{Terms}{误差项} & \tablelang{Group}{任务组} \\
\midrule
\tablelang{Harmonic grasp}{谐波抓取} & grasp & \tablelang{block}{物块} & \tablelang{table}{桌面} & \fillin{X} & -- & $v$, $r\theta$, $e_{\rm dyn}$ & R \\
\tablelang{Conveyor grasp}{传送带抓取} & t1\_belt & \tablelang{part}{零件} & \tablelang{belt}{传送带} & \fillin{X} & -- & $v$ & C \\
\tablelang{Translating grasp}{平移抓取} & t1\_xyz & \tablelang{part}{零件} & \tablelang{table}{桌面} & \fillin{X} & -- & $v$, $e_{\rm dyn}$ & -- \\
\tablelang{Rotating grasp}{旋转抓取} & t1\_so3 & \tablelang{part ($\omega\le1.2$\,rad/s)}{零件（$\omega\le1.2$\,rad/s）} & \tablelang{table}{桌面} & \fillin{X} & -- & $v$, $r\theta$, $e_{\rm dyn}$ & R \\
\tablelang{Occluded grasp}{遮挡抓取} & t1\_occ & \tablelang{part}{零件} & \tablelang{occluder canopy}{遮挡罩} & \fillin{X} & \fillin{X} & $v$, $m$, occ. & C \\
\tablelang{Moving receptacle}{移动容器} & t2\_trans & \tablelang{receptacle}{容器} & \tablelang{table}{桌面} & \fillin{X} & -- & $v$ & -- \\
\tablelang{Rotating receptacle}{旋转容器} & t2\_se3 & \tablelang{T-shaped receptacle}{T 形容器} & \tablelang{table}{桌面} & \fillin{X} & -- & $v$, $r\theta$ & R \\
\tablelang{Wire passage}{导线穿越} & t3\_phase & \tablelang{wire loop (entry)}{导线环（入口）} & \tablelang{rotating wire}{旋转导线} & 30 & 26.4$^\dagger$ & $v$, $r\theta$, $m$ & R, C \\
\tablelang{Wire passage, full}{导线穿越（完整）} & t3\_full & \tablelang{wire loop (entry)}{导线环（入口）} & \tablelang{rotating wire}{旋转导线} & 30 & 26.4$^\dagger$ & $v$, $r\theta$, $m$ & R, C \\
\tablelang{Ball capture}{小球捕获} & t4\_sphere & \tablelang{ball}{小球} & \tablelang{tray, ramp}{托盘、斜坡} & \fillin{X} & \fillin{X} & $v$, $e_{\rm dyn}$, $m$ & C \\
\tablelang{Ball capture, table}{小球捕获（桌面）} & t4\_sphere\_tabletop & \tablelang{ball}{小球} & \tablelang{table, ramp}{桌面、斜坡} & \fillin{X} & -- & $v$, $e_{\rm dyn}$ & C \\
\tablelang{Slider capture}{滑块捕获} & t4\_slider & \tablelang{slider}{滑块} & \tablelang{table, ramp}{桌面、斜坡} & \fillin{X} & -- & $v$, $e_{\rm dyn}$ & C \\
\tablelang{Can capture}{罐体捕获} & t4\_can & \tablelang{can}{罐体} & \tablelang{table, ramp}{桌面、斜坡} & \fillin{X} & -- & $v$, $r\theta$, $e_{\rm dyn}$ & R, C \\
\tablelang{Post-impact grasp}{撞击后抓取} & t5\_replan & \tablelang{part}{零件} & \tablelang{striker}{撞击器} & \fillin{X} & \fillin{X} & $v$, $r\theta$, $e_{\rm dyn}$, $m$, delay & R, C \\
\bottomrule
\end{tabularx}
\end{table*}


\subsection{Scenarios, Demonstrations, and Splits}
\label{app:splits}
\begin{table*}[!t]
\centering\fontsize{9}{11}\selectfont
\setlength{\tabcolsep}{4pt}
\renewcommand{\arraystretch}{1.12}
\caption{Historical validation proposal ranges. Length, linear velocity, angle,
and angular velocity use m, m/s, rad, and rad/s unless marked otherwise.
Positions are in world coordinates except the stated ramp-local offsets.
Continuous intervals are uniform proposals; discrete choices are equiprobable
before feasibility checks. Evaluation replays the stored accepted scenarios,
whose distribution need not remain uniform. Harmonic centers are not initial positions.}
\label{tab:task_sampling}
\begin{tabularx}{\textwidth}{@{}p{0.18\textwidth}p{0.34\textwidth}X@{}}
\toprule
Task & Initial-state proposal & Motion or event parameters \\
\midrule
grasp & Harmonic center $c_x\in[-.08,.08]$, $c_y\in[-.26,-.19]$, $z=.822$;
yaw $[-1.1,1.1]$.
& $v_y\in[.025,.125]$, $A_x\in[.015,.035]$, $f_x\in[.3,.7]$\,Hz;
$\omega_z\in[-1.2,1.2]$. \\
t1\_belt & Target $x\in[-.08,.08]$, $y\in[-.26,-.18]$, $z=.822$;
identity orientation, zero initial twist.
& Belt $v_y\in[.04,.10]$. \\
t1\_occ & Same target proposal as Belt; canopy at $(.322,-.29,1.20)$.
& Same belt range; clearance parameter .20; physical occlusion. \\
t1\_xyz & Harmonic center $c_x\in[-.025,.025]$, $c_y\in[-.24,-.20]$, $z=.818$;
yaw $[-.35,.35]$.
& $v_y\in[.04,.09]$, $A_x\in[0,.03]$, $f_x\in[.25,.65]$\,Hz; $\omega_z=0$. \\
t1\_so3 & Same as XYZ.
& Same translation; $\omega_z\in[-1.2,-.4]\cup[.4,1.2]$, equal sign probabilities. \\
t2\_trans & Goal $x\in[-.13,-.08]$, $y\in[-.18,-.12]$, $z=.825$; fixed orientation.
& Goal $v_y\in[.03,.06]$; no rotation. \\
t2\_se3 & Same goal $x,y$ range, $z=.804$; yaw $[-.3,.3]$.
& Goal $v_y\in[.03,.06]$;
$\omega_z\in[-.25,-.08]\cup[.08,.25]$, equal sign probabilities. \\
t3\_phase & Wand $x,y\in[-.03,.03]$, $z=.870$, yaw $[0,2\pi]$;
wire phase $[-.25,.35]$.
& Wire center $(.10,-.02,.94)$; world-X rotation
$\omega_x=-\{2.5,3,\ldots,7.5\}^{\circ}$/s (11 equal-probability values). \\
t3\_full & Wand at $(-.05,-.02,1.04)$ with identity orientation;
same wire phase range.
& Same wire rotation; cradle and release requirement enabled. \\
t4\_sphere & Ramp-local offsets: lateral $[-.015,.015]$, longitudinal/normal
$(-.145,.030)$; lip $(.06,-.02,1.10)$.
& Gravity-driven from rest; gate release at tick 0; slope $[5,15]^\circ$;
target/ramp sliding friction $[.7,.9]$; unsupported capture. \\
t4\_sphere\_tabletop & Lateral $[-.015,.015]$, longitudinal/normal
$(-.124,.034)$; lip $(0,-.06,.808)$.
& Same slope and release; target/ramp friction $[.7,.9]$, table friction .35. \\
t4\_slider & Lateral $[-.015,.015]$, longitudinal/normal $(-.129,.024)$;
supported ramp.
& Same slope and release; target/ramp friction $[.025,.045]$,
table friction $[.012,.018]$. \\
t4\_can & Lateral $[-.010,.010]$, ramp-axis yaw $[-4,4]^\circ$;
longitudinal position follows gate geometry and oriented can extent.
& Same slope and release; target/ramp friction $[.25,.45]$,
table friction $[.09,.15]$. \\
t5\_replan & $x\in[-.005,.005]$, $y=-.29$, $z=.812$;
yaw $[-.08,.08]$; zero initial twist.
& Belt $v_y\in[.12,.14]$; equal-probability left/right striker triggered at
$y=-.13$; stroke .098, pulse .35\,s, actuator scale $[1,1.1]$,
force cap 4\,N; command delay 50\,ms. \\
\bottomrule
\end{tabularx}
\end{table*}

\paragraph{Scenarios.} Table~\ref{tab:task_sampling} lists the scenario ranges. We
sample continuous ranges uniformly and discrete choices with equal probability, keep
the scenes that pass physical-feasibility checks, and replay the stored accepted
scenes at evaluation.

\paragraph{Demonstrations.} The expert is a task-specific closed-loop state machine
with privileged simulator state: the target's current pose and linear and angular
velocity, the robot state, the task geometry, and true contact, grasp, and stage
events. At every control period it, for example, extrapolates the target over a short
horizon and tracks, descends, closes, and lifts in the grasping tasks, predicts the
interception point of a rolling or flying object under gravity in the capture tasks,
and follows the known wire path with phase compensation from the current angle and
angular velocity in the passage tasks. Its commands are clipped to 0.05\,m and
0.5\,rad per axis. A demonstration is accepted if its scene is feasible, the expert
succeeds, an independent replay of the recorded actions reproduces the episode, and
its path, action changes, and jerk pass numerical quality checks. The training pool
of 11{,}200 demonstrations, 800 per task, is drawn from the accepted ones. We record
them at 20\,Hz with camera images, depth, robot states, issued commands, and simulator
ground truth for target poses and scene geometry; the learned policies receive only
the inputs of Sec.~\mainref{sec:method}.

\paragraph{Splits.} We select configurations on 1{,}400 development scenes (100 per
task). All \bench results use a separate held-out set of 1{,}400 scenes, generated with
the same ranges and new source seeds that appear in no training, validation, or
development episode; the first 100 feasible scenes in increasing seed order form each
task's set, which all methods share.

\subsection{Real-Time Evaluation}
\label{app:sweeps}
\paragraph{Executor.} While a method computes, physics advances by its latency and the
executor continues the previous chunk. All methods share the executor of
Sec.~\mainref{sec:inference} on the 50\,ms control clock, with queries on its period
boundaries; only the number $H_m$ of commands in a returned chunk differs
(Table~\ref{tab:baselines}). The total latency $L_t$ of query $t$ is its inference
latency under the protocol plus any added latency. Its chunk takes over at the
boundary $k_t\Delta_a$ after capture, with $k_t=\lceil L_t/\Delta_a\rceil$, from
command $k_t$, the one meant for that boundary; its earlier commands have expired and
are discarded, and a chunk with $k_t\ge H_m$ is dropped. The next query starts at that
boundary, or one control period later when $L_t=0$. Before the first chunk of an
episode arrives, and whenever a chunk runs out, the controller holds the current arm
target and gripper command while physics advances; no command is repeated. Command $j$
executes at $\tapp_{t,j}$ of Eq.~\mainref{eq:apptime}; the post-impact grasp applies
each command 50\,ms after it is issued ($d_{\rm task}=50$\,ms), and all other tasks
have $d_{\rm task}=0$.

\paragraph{Protocols.}
% Protocol definitions for Appendix B.3 (2026-10-25); results are in tables/supp_protocols.tex and supp_latency.tex.
\begin{table}[t]
\centering\footnotesize
\caption{\tablelang{\textbf{Evaluation protocols.} Latency: the time that physics advances while a method computes.}{\textbf{评测协议。}延迟：方法计算期间物理推进的时间。}}
\label{tab:protocol_defs}
\setlength{\tabcolsep}{3pt}
\begin{tabularx}{\columnwidth}{@{}l>{\raggedright\arraybackslash}Xl@{}}
\toprule
\tablelang{Protocol}{协议} & \tablelang{Latency}{延迟} & \tablelang{Executor}{执行器} \\
\midrule
\tablelang{Real time}{实时} & \tablelang{native, drawn per query}{原生延迟，每次查询抽取} & \tablelang{shared}{共享} \\
\tablelang{Added latency}{附加延迟} & \tablelang{native + 50--200\,ms}{原生延迟 + 50--200\,ms} & \tablelang{shared}{共享} \\
\tablelang{Paused}{暂停} & \tablelang{none (physics waits)}{无（物理等待）} & \tablelang{shared}{共享} \\
\tablelang{Equal latency}{等延迟} & \tablelang{p95 latency of \ours}{\ours 的 p95 延迟} & \tablelang{shared}{共享} \\
\tablelang{Original executor}{原执行方式} & \tablelang{native, drawn per query}{原生延迟，每次查询抽取} & \tablelang{method's own}{方法自身} \\
\bottomrule
\end{tabularx}
\end{table}

Table~\ref{tab:protocol_defs} defines the protocols. The native latency is drawn per
query from the method's latency distribution, measured on one \cfgGPU{}. Rollouts run on
\cfgGPU{} or \cfgEvalGPU{}; since physics advances only by the replayed latency, the
rollout hardware does not affect the results. The added-latency sweep of
Table~\ref{tab:supp_latency} adds 0, 50, 100, 150, or 200\,ms after inference under the
real-time protocol, and physics advances during it.

\paragraph{Latency measurement.}
\label{app:timing}
We time all methods on one \cfgGPU{} (32\,GB), the GPU of our robot, with exclusive use
and batch size 1. Each method runs three rounds of thirty warm-up and 300 measured
calls over three observation sequences of eight frames, with CUDA synchronization at
call boundaries and a rotated method order, in the precision of its release; \ours uses
FP32. A timed query starts from the newest observation, with the history already
maintained, and ends when its first command is available. It includes all computation
that needs the newest frame: preprocessing, image encoding, the online flow updates and
motion descriptors of PUMA and DynamicWAM, sampling, and action decoding, and on the
robot also tracking. For \ours, we time the forecast heads separately, although
deployment skips them. This time-to-action is reported without rounding; the
handover time $\ell_t=k_t\Delta_a$ rounds it up to the control period, and the action
age of a command is its execution time minus the capture time. A query misses its
deadline when its time-to-action exceeds 50\,ms. Appendix~\ref{app:latency_results}
reports the measured values.

\subsection{Metrics and Statistics}
\label{app:metrics}
The task-balanced success rate is
\begin{equation}
\mathrm{SR}(m)=\frac{100}{14}\sum_{k=1}^{14}\frac1{S_m n_k}\sum_{s=1}^{S_m}\sum_{i=1}^{n_k}y_{mski},
\end{equation}
for binary outcomes $y_{mski}$ of method $m$, trained with seed $s$, on scene $i$ of
task $k$, with $S_m=3$ training seeds for our variants and $S_m=1$ for the
baselines. The
\emph{position error at first contact} $e_p$ is the distance of the end effector from
the expert's contact frame on the target at the first contact. The two placement
tasks start with the object held, so their contact is the first contact of the held
object with the receptacle, and their contact frame is the end-effector pose implied
by the expert's placement target. The phase-only passage (\texttt{t3\_phase})
succeeds without any contact, so $e_p$ averages over the other 13 tasks; we report it
on episodes that make contact, and failures without contact stay in the SR
denominator.

\paragraph{Training seeds and intervals.}
\label{app:statistics}
The full model and the five ablation controls use training seeds 0, 1, and 2, and each
baseline uses seed 0. We report each training seed and the cross-seed mean and
standard deviation (Appendix~\ref{app:ablations}), and for the ablations also the
comparison without the phase-only wire passage (\texttt{t3\_phase}). For paired 95\%
intervals, we average the seeds per scene and bootstrap held-out scenes within each
task, jointly across methods. These intervals quantify scene uncertainty for the
trained weights, and we treat differences whose interval lies within $\pm2$ points as
ties.

\subsection{Baselines and Ablation Controls}
\label{app:baselines}
% Protocol table: Dynamic-14 baseline adaptations from BASE1 and the v3 endpoint records.
% PUMA uses the resolved 40k training configuration and parameter summary in
% BASE1/results/puma_dynamic14_20261101/h100/training-final.zip.
\begin{table}[!htbp]
\centering\footnotesize
\caption{\tablelang{\textbf{Baseline training.} The training budget gives optimizer updates and global batch size.
$H$ is the number of commands returned per query. The parameter counts include trainable parameters only.
DOM ckpt. denotes the released DynamicVLA checkpoint trained on DOM.}{\textbf{基线训练。}训练量以优化器更新次数和全局 batch 大小表示。
$H$ 为每次查询返回的指令数。参数量只计入可训练参数。DOM 检查点指 DynamicVLA 发布的、在 DOM 上训练的检查点。}}
\label{tab:baselines}
\setlength{\tabcolsep}{2.6pt}
\begin{tabular}{@{}llccc@{}}
\toprule
\tablelang{Method}{方法} & \tablelang{Initialization}{初始化} & \tablelang{Updates$\times$batch}{更新数$\times$batch} & $H$ & \tablelang{Params}{参数量} \\
\midrule
ACT \cite{act} & \tablelang{scratch}{从头训练} & $50{,}000\times32$ & 100 & 83.9M \\
DP \cite{diffusionpolicy} & \tablelang{scratch}{从头训练} & $50{,}000\times32$ & 15 & 288.7M \\
$\pi_{0.5}$ \cite{pi05} & \texttt{pi05\_base} & $30{,}000\times256$ & 10 & 3,353.4M \\
$\pi_{0.5}$+SE(3) & \texttt{pi05\_base} & $30{,}000\times256$ & 10 & 3,357.7M \\
DynamicVLA \cite{dynamicvla} & \tablelang{DOM ckpt.}{DOM 检查点} & $50{,}000\times44$ & 20 & 99.1M \\
PUMA \cite{domino} & \tablelang{official}{官方} & $40{,}000\times64$ & 16 & 4,516.5M \\
OpenWAM-$\alpha$ \cite{openwam_alpha} & \tablelang{official}{官方} & $34{,}490\times192$ & 32 & 6,021.2M \\
DynamicWAM \cite{dynamicwam} & \tablelang{official}{官方} & $40{,}000\times128$ & 16 & 988.8M \\
DynamicWAM+SE(3) & \tablelang{official}{官方} & $40{,}000\times128$ & 16 & 989.4M \\
\midrule
\ours & \tablelang{scratch}{从头训练} & $\cfgUpdates\times512$ & \cfgChunk & \cfgTrainableParams \\
\bottomrule
\end{tabular}
\end{table}

\paragraph{Common conditions.} Each baseline keeps its official architecture, chunk
length, initialization, and update budget, trains once with seed 0 on the same
11{,}200 demonstrations, whose commands we map to its action space through the same
Cartesian controller, and is configured on development scenes. All methods share the
agent-view camera; Table~\ref{tab:baselines} lists their further inputs, the handling
of missing wrist views, and method-specific settings.

\paragraph{SE(3)-augmented baselines.} $\pi_{0.5}$+SE(3) and DynamicWAM+SE(3)
receive the eight target-state tokens of Appendix~\ref{app:tokens} from the same
pose source, with the same validity flags and masking. A two-layer MLP maps the
tokens to the model's width. For $\pi_{0.5}$, we append them to the prefix after the
image tokens, where the action expert attends to them; for DynamicWAM, they join its
kinematic descriptors. The output layer of the MLP is zero-initialized, and training
follows the base model's recipe.

\paragraph{Ablation controls.}
\label{app:ablation_defs}
The five controls of Sec.~\mainref{sec:principles} change one component each of the
configuration of Table~\ref{tab:hyper} and keep the data, the recipe, the update
budget, the executor, and the training seeds 0, 1, and 2. \emph{No target history}
removes slots 0--6 of the target-state tokens, which carry the past poses with their
velocity and acceleration features, and keeps slot 7 with the current pose
(Appendix~\ref{app:tokens}). \emph{No future supervision} removes $\mathcal L_P$,
$\mathcal L_F$, and the three future appearance losses with their heads, and keeps
the reconstruction of the current target region. \emph{No scene-flow supervision}
removes $\mathcal L_F$ and the flow readout, and \emph{no dedicated SE(3)} removes
$\mathcal L_P$ and the pose head; both keep the target-state tokens and all
appearance targets. The \emph{uniform-fidelity forecast} replaces $\mathcal L_P$ and
$\mathcal L_F$ by \fillin{a forecast of future latents of a frozen video VAE of the
whole frame at the same horizons, its readout, and its loss weight}. Since every
forecast head is skipped at deployment, only the no-history control changes the
inputs at inference.
